%%%%%%%%%%%%%%%%%%%%%%%%
%%%%%%% PACKAGES %%%%%%%
%%%%%%%%%%%%%%%%%%%%%%%%

%\usepackage{babel}															% Languages, necessary letters.
\usepackage{braket}															% Adds Dirac bra-kets.
\usepackage{lipsum}															% Adds dummy text.
\usepackage{accents}														% Set symbols above characters; for vector fix.
\usepackage{titlecaps}													% Commands for Title Case.
\usepackage{gensymb}														% Symbols for some physics units.
\usepackage{unicode-math}												% Math fonts with unicode symbols.
\usepackage{slashed}														% Adds a slash for Feynman notation.
\usepackage{titling}														% Adjust title layout.
\usepackage{stackengine}												% Stack symbols; for \decileft (below)
\usepackage{xfrac}															% More fractions, e.g. slanted fractions.
\usepackage{multirow}														% Extends table environment.
% \usepackage{mathexam}														%Exam formatting
\usepackage{slashed}														% Feynman slash notation.
\usepackage{xcolor}															% More colors.
\usepackage{siunitx}														% Units & scientific notation.
\usepackage{tikz}																% Graphics.
\usepackage{yfonts}															% German fonts, e.g. Fraktur.
\usepackage{lettrine}														% Adds big fancy script letters.
\usepackage{adjustbox}													% Customization of boxing.
\usepackage{aurical}														% Adds a calligraphic font
\usepackage{anyfontsize}												% Allows for using *any* font size.
\usepackage{graphicx}														% Additional graphics options.
\usepackage{tikz-feynman}												% Feynman diagrams.
\usepackage[framemethod=TikZ]{mdframed}					% Customizable boxing
\usepackage[object=vectorian]{pgfornament}			% Page ornaments.
\usepackage[notquote]{hanging}		%Hanging paragraphs; option prevents choking on apostrophes.


%%%%%%%%%%%%%%%%%%%%%%%%%%%%%%%%%
%%%%%%% PACKAGE LIBRARIES %%%%%%%
%%%%%%%%%%%%%%%%%%%%%%%%%%%%%%%%%
\usetikzlibrary{
	backgrounds,
	arrows.meta,
	shapes.misc,
	bending,
	positioning,
	quotes,
	angles,
	graphs.standard,
	decorations.pathmorphing,
	decorations.markings,
	calc,
	spath3,
	intersections
}

\usepgflibrary{plotmarks}

% Prevents a margin error with the geometry package
\geometry{nomarginpar}

%%%%%%%%%%%%%%%%%%%%%%%%%%%%%%%%
%%%%%%% PACKAGE SETTINGS %%%%%%%
%%%%%%%%%%%%%%%%%%%%%%%%%%%%%%%%

% Fix for yinit/yfonts; allows for large script letters.
\def\initdefault{yinit}

% Sets the title for exams.
% \ExamClass{Phys 105}
% \ExamName{Practice Midterm 2}
% \ExamHead{Spring 2026}


%%%%%%%%%%%%%%%%%%%%%%
%%%%%%% LAYOUT %%%%%%%
%%%%%%%%%%%%%%%%%%%%%%
	
% Remove space before title
% \setlength{\droptitle}{-15em} 
% \posttitle{\par\end{center}}


%%%%%%%%%%%%%%%%%%%%%%%%%%%%%%%%%%
%%%%%%% DOC COLORS & FONTS %%%%%%%
%%%%%%%%%%%%%%%%%%%%%%%%%%%%%%%%%%


%%%%%%%% Produces dark mode pdfs.
%\pagecolor{black}
%\color{white}

% Set main math font as KPMath Sans.
% This will only be used as a fallback,
% if the fonts after are missing characters.
\setmathfont{KPMath-Sans}
% Use XITS Math font for some specific symbols.
\setmathfont{XITS Math}[range={\acwunderarcarrow,\uprightcurvearrow}]
% For the rest, default to the OpenDyslexic font.
\setmathfont{OpenDyslexic Italic}[range=it]
\setmathfont{OpenDyslexic}[range=up]
\setmathfont{OpenDyslexic}[range={"01-"02,"04-"27,"30-"7B,"7D-"FF,"2B,"2D,"2E,"B7}]
\setmathfont{KPMath-Sans}[range={"03D1,"1D751}]



%%%%%%%%%%%%%%%%%%%%%%%%%%%%%
%%%%%%% CUSTOM STYLES %%%%%%%
%%%%%%%%%%%%%%%%%%%%%%%%%%%%%

% Custom boxing style.
\mdfdefinestyle{important}{%
  leftmargin=3cm,%
	rightmargin=3cm,%
	linewidth=2pt,%
	roundcorner=10pt,%
	backgroundcolor=gray!30
}

\newcommand{\sectionline}{
  \noindent
  \begin{center}{
  	\color{DarkViolet}
    \resizebox{0.5\linewidth}{1ex}{
    	\begin{tikzpicture}
    		\node (C) at (0,0) {};
    		\node (D) at (9,0) {};
    		\path (C) to [ornament=85] (D);
    	\end{tikzpicture}
    }
  }\end{center}
}

%% A macro with two arguments to change ornaments and colors easily
%% Syntax -- \sectionlinetwo{<color>}{<ornament>}
\newcommand{\sectionlinetwo}[2]{
  \nointerlineskip \vspace{.5\baselineskip}\hspace{\fill}{
  	\color{#1}
    \resizebox{0.5\linewidth}{2ex}{
    	\begin{tikzpicture}
    		\node (C) at (0,0) {};
    		\node (D) at (9,0) {};
    		\path (C) to [ornament=#2] (D);
    	\end{tikzpicture}
    }
    \hspace{\fill}
    \par\nointerlineskip \vspace{.5\baselineskip}
  }
}

% Set style of lettrine to the yint font.
\renewcommand{\LettrineFontHook}{\usefont{U}{yinit}{m}{n}}
% Adjust default size of lettrine
\renewcommand*{\DefaultLoversize}{0.1}
% Adjust default number of lines that the lettrine takes up
\setcounter{DefaultLines}{4}
% Adjust space between lettrine and text immediately after
\AtBeginDocument{\setlength{\DefaultFindent}{0.5em}}
% Adjust indent of the remaining lines
\setlength{\DefaultNindent}{0pt}
% Shift lettrine up
\renewcommand{\DefaultLraise}{-0.01}


%%%%%%%%%%%%%%%%%%%%%%%%%%%
%%%%%%% TIKZ MACROS %%%%%%%
%%%%%%%%%%%%%%%%%%%%%%%%%%%

%%%%%%%% Shapes for nodes.

% Custom dots
\tikzset{dots/.style={shape=circle,fill=black,scale=0.3}}
% Custom cross
\tikzset{cross/.style={
	cross out, draw=black,
	minimum size = 0.25em,
	inner sep=0pt, outer sep=0pt
}}
% Custom triangle
\tikzset{triangle/.style={
	regular polygon,
	regular polygon sides=3
}}
% Custom diamond
\tikzset{diamond/.style={
	regular polygon,
	regular polygon sides=4,
	shape border rotate=45
}}

% Draws a box
% Syntax - \boxFloor{bottom center pos}{size}{draw options}
\newcommand{\boxFloor}[3]{
	\draw[#3] ($(#1) - (#2/2,0)$) rectangle ($(#1) + (#2/2,#2)$)
}


% Draws a box
% Syntax - \boxFloor{center pos}{size}{draw options}
\newcommand{\boxCenter}[3]{
	\draw[#3] ($(#1) - (#2/2,#2/2)$) rectangle ($(#1) + (#2/2,#2/2)$)
}

%%%%%%%% For geometry & coordinate planes

% Custom angle marker
% Syntax -- \anglabel{<vertex coord>}{<start angle>}{<end angle>}{<arc radius>}
\NewDocumentCommand{\anglabel}{O{} m m m m}{
	\draw[#1]($#2 + (#3:#5)$) arc (#3:#4:#5)
}

% Custom dashes
\tikzset{dashV/.style={midway, shape=rectangle, inner xsep=-10pt, fill}}
\tikzset{dashH/.style={midway, shape=rectangle, inner ysep=-10pt, fill}}
% Syntax -- \draw... node[dash[<rotation angle of dash>],...]
\tikzset{dash[1]/.style={midway, shape=rectangle, inner ysep=-10pt, fill, shape border rotate=#1}}

% Draws a coordinate plane
% Assumes that axes are of equal size
% Syntax -- \coordAxesSq{<origin coords>}{<size of +axis>}
\newcommand{\coordAxesReg}[2]{

	\coordinate(origin) at (#1);
	\coordinate(minusX) at ($(origin)+(-#2,0)$);
	\coordinate(plusX) at ($(origin)+(#2,0)$);
	\coordinate(minusY) at ($(origin)+(0,-#2)$);
	\coordinate(plusY) at ($(origin)+(0,#2)$);

	\draw[ultra thick, -{Latex}] (minusX) -- (plusX)node[right]{$x$};
	\draw[ultra thick, -{Latex}] (minusY) -- (plusY)node[above]{$y$};

}

\tikzset{coordAxesSq/.pic = {
	\draw[-{Latex}](-1,0) -- (1,0)node[right]{$x$};
	\draw[-{Latex}](0,-1) -- (0,1)node[above]{$y$};
}}

% Draws a coordinate plane
% Syntax -- \coordplaneReg{<origin coords>}{< -x bound>}{< +x bound>}{< -y bound>}{< +y bound>}
\newcommand{\coordAxes}[5]{

	\coordinate(origin) at (#1);
	\coordinate(minusX) at ($(origin)+(#2,0)$);
	\coordinate(plusX) at ($(origin)+(#3,0)$);
	\coordinate(minusY) at ($(origin)+(0,#4)$);
	\coordinate(plusY) at ($(origin)+(0,#5)$);

	\draw[ultra thick, -{Latex}] (minusX) -- (plusX)node[right]{$x$};
	\draw[ultra thick, -{Latex}] (minusY) -- (plusY)node[above]{$y$};

}

% Draws a square grid of a specified size, centered at the origin
% Syntax -- \coordplane{<origin coords>}{<size of +axis}
\newcommand{\coordGridSq}[2]{
	\foreach \x in {-#2,...,#2}{
		\draw[line width=0.25pt,shift={#1}](\x,-#2)--(\x,#2);
	}		
	\foreach \y in {-#2,...,#2}{
		\draw[line width=0.25pt,shift={#1}](-#2,\y)--(#2,\y);
	}
}

% Draws a square grid of a specified size, centered at the origin
% Syntax -- \coordplane{<origin coords>}{< -x bound>}{< +x bound>}{< -y bound>}{< +y bound>}
\newcommand{\coordGrid}[5]{
	\foreach \x in {#2,...,#3}{
		\draw[line width=0.25pt,shift={#1}](\x,#4)--(\x,#5);
	}		
	\foreach \y in {#4,...,#5}{
		\draw[line width=0.25pt,shift={#1}](#2,\y)--(#3,\y);
	}
}

% Draws a square dot-grid of a specified size, centered at the origin
% Syntax -- \coordplane{<origin coords>}{<size of +axis}
\newcommand{\coordGridDotsSq}[2]{
	\foreach \x in {-#2,...,#2}{	
		\foreach \y in {-#2,...,#2}{
			\draw[shift={#1}](\x,\y)node[dots]{};
		}
	}
}

% Draws a square dot-grid of a specified size, centered at the origin
% Syntax -- \coordplane{<origin coords>}{< -x bound>}{< +x bound>}{< -y bound>}{< +y bound>}
\newcommand{\coordGridDots}[5]{
	\foreach \x in {#2,...,#3}{	
		\foreach \y in {#4,...,#5}{
			\draw[shift={#1}]node[dots](\x,\y){};
		}
	}
}

%%%%%%%% Special line types.

% Draws a "rough" line.
% Syntax -- \drawRough{startcoord}{endcoord}{roughness}
\NewDocumentCommand{\drawRough}{m m m O{}}{
	\draw[#4,
		postaction=decorate,
		decoration={
			markings,
			mark=between positions 1cm and 3cm step 0.5cm with {
				\node[
					transform shape,
					inner sep=1pt
				](
					mark-\pgfkeysvalueof
						{/pgf/decoration/mark info/sequence number}
				){};
			}
		}
	]{
		decorate[
			decoration={
				random steps,
				segment length=1.5mm,
				amplitude=#3 %Roughness is amplitude
			}
		]{(#1)--(#2)}
	}
}

% Draws a "fat" arrow

\tikzset{fatArrow/.style={
	-{Triangle[
		width=\the\dimexpr1.8\pgflinewidth,
		length=\the\dimexpr0.8\pgflinewidth
	]}
}}

\NewDocumentCommand{\thickArrow}{O{} m m m}{
	\coordinate(thickArrow) at ($(#3)!1.6*#4cm!(#2)$);
	\draw[#1]
		($(thickArrow)!#4cm!-90:(#2)$)--
		($(#2)!#4cm!90:(thickArrow)$)--
		($(#2)!#4cm!-90:(thickArrow)$)--
		($(thickArrow)!#4cm!90:(#2)$)--
		($(thickArrow)!1.4*#4cm!90:(#2)$)--
		($(thickArrow)!1.6*#4cm!180:(#2)$)--
		($(thickArrow)!1.4*#4cm!-90:(#2)$)--
		cycle;
}

%%%%%%%%%%%%%% Pre-defined graphics objects.

% Stick figure

\tikzset{stickFig/.pic ={
	\draw(0,0.25)--(0,-0.25);
	\draw(0,0.375)circle(0.125);
	\draw(0.25,-0.5)--(0,-0.25)--(-0.25,-0.5);
	\draw(0.25,0)--(-0.25,0);
}}

% Apple
\tikzset{apple/.pic={
	\fill [brown] (-1/8,0) 
		arc (180:120:1 and 3/2) coordinate [pos=3/5] (@)-- ++(1/6,-1/7) 
		arc (120:180:5/4 and 3/2) -- cycle;
	\fill [red] (0,-9/10) 
		.. controls ++(180:1/8) and ++(  0:1/4) .. (-1/3,  -1)
		.. controls ++(180:1/3) and ++(270:1/2) .. (  -1,   0)
		.. controls ++( 90:1/3) and ++(180:1/3) .. (-1/2, 3/4)
		.. controls ++(  0:1/8) and ++(135:1/8) .. (   0, 4/7)
		.. controls ++( 45:1/8) and ++(180:1/8) .. ( 1/2, 3/4)
		.. controls ++(  0:1/3) and ++( 90:1/3) .. (   1,   0)
		.. controls ++(270:1/2) and ++(  0:1/3) .. ( 1/3,  -1)
		.. controls ++(180:1/4) and ++(  0:1/8) .. cycle;
	\fill [red] (0, 4/7)
		.. controls ++( 45:1/8) and ++(180:1/8) .. ( 1/2, 3/4)
		.. controls ++(  0:1/3) and ++( 90:1/3) .. (   1,   0)
		.. controls ++(270:1/2) and ++(  0:1/3) .. ( 1/3,  -1)
		.. controls ++(180:1/4) and ++(  0:1/8) .. (   0,-9/10);
	\fill [green, shift={(@)}, rotate=-30] 
		(0,0) arc (45:135:3/4 and 3/5) arc (225:315:3/4 and 3/5);
	\fill [green, shift={(@)}, rotate=-30] 
		(0,0) arc (315:225:3/4 and 3/5) -- cycle;
}}

% Car
\tikzset{car/.pic={
	\draw
		(6,-4.5)--++
		(0,0.5)--++
		(0.75,0)--++
		(45:0.5)--++
		(1,0)--++
		(-45:0.5)--++
		(0.75,0)--++
		(0,-0.75)--++
		(-0.5,0)arc
		(0:180:0.25)--++
		(-1,0)arc
		(0:180:0.25)--++
		(-0.705,0)--++
		(0,0.25);
	\draw(8.46,-4.75)circle(0.225);
	\draw(6.96,-4.75)circle(0.225);
}}

% Eye
\tikzset{eye/.pic={
	\draw (0,0)--(15:1.25);
	\draw (0,0)--(-15:1.25);
	\draw ($(0,0) + (-15:1)$) arc (-15:15:1);
	\draw ($(1.5,0) + (195:0.775)$) arc (195:165:0.775);
	\fill ($(1,0)+(-0.0625,0)$) ellipse (0.0625 and 0.125);
}}

% Ramp
% Syntax = \blockRamp{ramp right angle corner}{rampAngle}{rampLength}
\NewDocumentCommand{\ramp}{m m m}{
	\draw
		(#1)--++
		(0,{#3*sin(#2)})--++
		({#3*cos(#2)},{-#3*sin(#2)})--
		(#1);
	\anglabel{($(#1)+({#3*cos(#2)},0)$)}{180}{180-#2}{0.125*#3}
}

% Block on a ramp
% Syntax = \blockRamp{ramp right angle corner}{blockSize}{% from top}{rampAngle}{rampLength}
\NewDocumentCommand{\blockRamp}{O{} m m m m m}{
	\draw[#1]
		(#2)--++
		(0,{#6*sin(#5)})--++
		({#6*cos(#5)},{-#6*sin(#5)})--
		(#2);
	\draw
		($
			(#2)+
			(0,{#6*sin(#5)})+
			#4*#6*({cos(#5)},{-sin(#5)})+
			#3*({-cos(#5)},{sin(#5)})
		$)--++
		({#3*sin(#5)},{#3*cos(#5)})--++
		({#3*cos(#5)},{-#3*sin(#5)})--++
		({-#3*sin(#5)},{-#3*cos(#5)})--++
		({-#3*cos(#5)},{#3*sin(#5)});
	\anglabel{($(#2)+({#6*cos(#5)},0)$)}{180}{180-#5}{0.125*#6}
}


%%%%%%%%%%%%%%%%%%%%%
%%%%%%% FIXES %%%%%%%
%%%%%%%%%%%%%%%%%%%%%
	
% Fix primed vector spacing issue. Uses accents package.
\let\oldvec\vec
\renewcommand{\vec}[1]{\oldvec{#1}\mkern2mu\vphantom{#1}}

% Increase spacing in matrices
\renewcommand{\arraystretch}{1}


%%%%%%%%%%%%%%%%%%%%%%%%%%%%%%%
%%%%%%% CUSTOM COMMANDS %%%%%%%
%%%%%%%%%%%%%%%%%%%%%%%%%%%%%%%

%%%%%%%% Math additions

%Shortcut for scientific notation
\newcommand\e[1]{
    \times{10}^{#1}
}

% Used for showing "movement" of the decimal point in scientific notation.
% Requires a unicode math font (SEE ABOVE)
\newcommand\decileft[1]{%
    \kern-.4ex\stackunder[0.4pt]{$#1$}{$\color{red}\kern0.2ex\acwunderarcarrow$}
}

\newcommand\deciright[1]{%    <--- Decimal position to left
    \kern-.4ex\stackunder[0.4pt]{$#1$}{%
      \reflectbox{$\color{red}\kern-1.4ex\acwunderarcarrow$}
      }
}

% Operators/functions with Roman text, as with \cos.
\DeclareMathOperator{\tr}{tr}
\DeclareMathOperator{\arcsinh}{arcsinh}
\DeclareMathOperator{\arccosh}{arccosh}
\DeclareMathOperator{\arctanh}{arctanh}
\DeclareMathOperator{\arccsch}{arccsch}
\DeclareMathOperator{\arcsech}{arcsech}
\DeclareMathOperator{\arccoth}{arccoth}


%%%%%%%% Symbols

% Unicode math redefines vartheta. Define the version we want here:
\newcommand{\ctheta}{\symbol{"1D751}}

% Used for ""-bar characters.
\newcommand{\smash@bar}[4]{%
  \smash{\rlap{\raisebox{-#3\fontdimen5#10}{$\m@th#2\mkern#4mu\mathchar'26$}}}%
}

% Adds lambda-bar character.
\newcommand{\lambdabar}{{\mathchoice
  {\smash@bar\textfont\displaystyle{0.25}{1.2}\lambda}
  {\smash@bar\textfont\textstyle{0.25}{1.2}\lambda}
  {\smash@bar\scriptfont\scriptstyle{0.25}{1.2}\lambda}
  {\smash@bar\scriptscriptfont\scriptscriptstyle{0.25}{1.2}\lambda}
}}

% Adds d-bar character.
\newcommand{\dbar}{{\mathchoice
  {\smash@bar\textfont\displaystyle{0.1}{2.5}d}
  {\smash@bar\textfont\textstyle{0.1}{2.5}d}
  {\smash@bar\scriptfont\scriptstyle{0.1}{2.5}d}
  {\smash@bar\scriptscriptfont\scriptscriptstyle{0.1}{2.5}d}
}}

%%%%%%%% Figure/Table/Equation Numbering

% Repeat figures and tables without new numbering
\newcommand{\repeatcaption}[3]{%
  \captionof*{#1}{\titlecap{#1} \ref{#2} (revisited): #3}
}


% % Resets equation numbering per section
% \@addtoreset{equation}{section}
% % Adds section number to equation numbers.
% \numberwithin{equation}{section}

%set the "per" mode in siunitx to use the slash
\sisetup{per-mode = symbol}


%%%%%%%%% SI Unit Definitions, Settings, etc.

% Shortcut for units to a power
\newcommand{\p}[1]{\tothe{#1}}
% Shortcut for degrees symbol
\newcommand{\dg}{\degree}
% Shortcuts for common physics units
\newcommand{\vel}{\m\per\s}
\newcommand{\accel}{\m\per\s\tothe{2}}
\newcommand{\force}{\kg.\m\per\s\tothe{2}}
\newcommand{\energy}{\kg.\m\tothe{2}\per\s\tothe{2}}

% Define imperial units
% length
\DeclareSIUnit{\inch}{in}
\DeclareSIUnit{\ft}{ft}
\DeclareSIUnit{\yd}{yard}
\DeclareSIUnit{\mi}{mi}

% time
\DeclareSIUnit{\min}{min}
\DeclareSIUnit{\hr}{h}
\DeclareSIUnit{\day}{d}
\DeclareSIUnit{\wk}{week}
\DeclareSIUnit{\yr}{yr}

% mass
\DeclareSIUnit{\lb}{lb}
\DeclareSIUnit{\oz}{oz}

% volume
\DeclareSIUnit{\floz}{fl oz}
\DeclareSIUnit{\gal}{gal}

% temperature
\DeclareSIUnit{\F}{F}

% acceleration due to gravity
\newcommand{\grav}{\qty{9.8}{\m\per\s\tothe{2}}}